\scanpage{149}
% archive sheet number: 148
% L1–4
et $\widehat{\mathbb Z}.\dot\varepsilon_0$ (isom. au groupe des $\left(\begin{smallmatrix}1&\mu\\0&1\end{smallmatrix}\right)$) que l'action de $\widehat{\mathbb Z}^{*}$ sur $\dot\varepsilon_0$ est bien triviale -- ce qui était clair a priori -- donc l'action de $\widehat{\mathbb Z}^{*}/\pm1$ dans (53) s'exprime par un caractère
% L5
\[
\varepsilon:\widehat{\mathbb Z}^{*}/\pm1\longrightarrow\{\pm1\}\deleted{\unsure{\varepsilon_i}}
\tag*{(54)}
\]
% L6–9
qu'il faudrait expliciter, exprimant de quelle façon un élément arbitraire \deleted{\unsure{de}} $\dot\lambda\in\widehat{\mathbb Z}^{*}/\pm1$ opère sur $\{\underset{\substack{\Vert\\\dot\tau_0}}{\dot\tau_1},\dot\tau_\infty\}$, \deleted{-- rappelons-nous à ce propos que}
% L10
\[
\deleted{r_0=\begin{pmatrix}0&1\\1&0\end{pmatrix},\quad
 r_1=\begin{pmatrix}-1&1\\0&1\end{pmatrix},\quad
 r_\infty=\begin{pmatrix}-1&0\\0&1\end{pmatrix}.}\tag*{\deleted{(55)}}
\]
% L11–17
\deleted{Si $\varepsilon$ n'est pas le caractère trivial,}

\deleted{la structure connue de $\widehat{\mathbb Z}^{*}$ montre qu'il suffit de tester la valeur de $\varepsilon$ pour les $\mathbb Z/n\mathbb Z$}

\deleted{A vrai dire, il est clair a priori que \deleted{la valeur} $\varepsilon$ est trivial sur les $\mathbb Z_\ell^{*}\simeq\mathbb Z_\ell\times\mathbb F_\ell^{*}$ avec $\ell$ nb premiers impairs [car il est trivial, sur $\mathbb Z_\ell$ la composante $\mathbb Z_\ell$, ce qui montre qu'il suffit de regarder}
% L18–19
\deleted{Écrivons}
\[
\deleted{\widehat{\mathbb Z}^{*}\simeq\widehat{\mathbb Z}_2^{*}\times
\underbrace{\widehat{\mathbb Z}^{*\,\prime}}_{\displaystyle
\prod_{\substack{\ell\text{ premier}\\\text{impair}}}\mathbb Z_\ell^{*}}\, ,}
\]
% L20–24
\deleted{il est clair que si $\varepsilon\ne1$, alors $\varepsilon$ ne provient \uline{pas} d'un caractère sur le facteur $\widehat{\mathbb Z}^{*\,\prime}$ (i.e. il ne provient pas d'un caractère sur $(\mathbb Z/n\mathbb Z)^{*}$ avec $n$ \deleted{premier} impair), donc dans ce cas la composante du caractère suivant $\widehat{\mathbb Z}_2^{*}$}
